% Principal axes of a uniform ellipsoid and the meaning of the principal moments.
% Semi-axes a > b > c along x1, x2, x3, so that I1 < I2 < I3 (I1 = mu (b^2+c^2)/5, etc.).
% Hidden parts of the principal sections are dashed; the silhouette is computed exactly
% from the view direction, so the drawing stays correct if the view angles are changed.
% Compile: xelatex principal_axes_tikz.tex ; pdftoppm -png -r 300 -singlefile principal_axes_tikz.pdf principal_axes
\documentclass[tikz,border=4pt]{standalone}
\usepackage{amsmath}
\usepackage[fontset=fandol]{ctex}
\usepackage{tikz-3dplot}
\usetikzlibrary{arrows.meta}

\begin{document}
\def\vth{70}\def\vph{-25}          % view angles
\tdplotsetmaincoords{\vth}{\vph}
\begin{tikzpicture}[tdplot_main_coords,>=Stealth,line cap=round,
                    every node/.style={font=\normalsize}]
  \def\a{2.6}\def\b{1.6}\def\c{1.05}
  % view direction v (towards the reader) and g = A^{1/2} v
  \pgfmathsetmacro{\vx}{sin(\vth)*sin(\vph)}
  \pgfmathsetmacro{\vy}{-sin(\vth)*cos(\vph)}
  \pgfmathsetmacro{\vz}{cos(\vth)}
  \pgfmathsetmacro{\gx}{\vx/\a}\pgfmathsetmacro{\gy}{\vy/\b}\pgfmathsetmacro{\gz}{\vz/\c}
  \pgfmathsetmacro{\gn}{sqrt(\gx*\gx+\gy*\gy+\gz*\gz)}
  \pgfmathsetmacro{\gx}{\gx/\gn}\pgfmathsetmacro{\gy}{\gy/\gn}\pgfmathsetmacro{\gz}{\gz/\gn}
  % e = g x z / |g x z|,  f = g x e
  \pgfmathsetmacro{\en}{sqrt(\gx*\gx+\gy*\gy)}
  \pgfmathsetmacro{\ex}{\gy/\en}\pgfmathsetmacro{\ey}{-\gx/\en}\pgfmathsetmacro{\ez}{0}
  \pgfmathsetmacro{\fx}{\gy*\ez-\gz*\ey}\pgfmathsetmacro{\fy}{\gz*\ex-\gx*\ez}\pgfmathsetmacro{\fz}{\gx*\ey-\gy*\ex}
  % silhouette and the three principal sections
  \fill[blue!7] plot[domain=0:360,samples=120,variable=\t]
     ({\a*(cos(\t)*\ex+sin(\t)*\fx)},{\b*(cos(\t)*\ey+sin(\t)*\fy)},{\c*(cos(\t)*\ez+sin(\t)*\fz)}) -- cycle;
  \draw[blue!70!black,thick] plot[domain=0:360,samples=120,variable=\t]
     ({\a*(cos(\t)*\ex+sin(\t)*\fx)},{\b*(cos(\t)*\ey+sin(\t)*\fy)},{\c*(cos(\t)*\ez+sin(\t)*\fz)}) -- cycle;
  % section in the x1x2 plane: visible for |t - t0| < 90
  \pgfmathsetmacro{\tA}{atan2(\vy/\b,\vx/\a)}
  \draw[blue!45] plot[domain=\tA-90:\tA+90,samples=60,variable=\t] ({\a*cos(\t)},{\b*sin(\t)},0);
  \draw[blue!45,dashed] plot[domain=\tA+90:\tA+270,samples=60,variable=\t] ({\a*cos(\t)},{\b*sin(\t)},0);
  % section in the x1x3 plane
  \pgfmathsetmacro{\tB}{atan2(\vz/\c,\vx/\a)}
  \draw[blue!45] plot[domain=\tB-90:\tB+90,samples=60,variable=\t] ({\a*cos(\t)},0,{\c*sin(\t)});
  \draw[blue!45,dashed] plot[domain=\tB+90:\tB+270,samples=60,variable=\t] ({\a*cos(\t)},0,{\c*sin(\t)});
  % section in the x2x3 plane
  \pgfmathsetmacro{\tC}{atan2(\vz/\c,\vy/\b)}
  \draw[blue!45] plot[domain=\tC-90:\tC+90,samples=60,variable=\t] (0,{\b*cos(\t)},{\c*sin(\t)});
  \draw[blue!45,dashed] plot[domain=\tC+90:\tC+270,samples=60,variable=\t] (0,{\b*cos(\t)},{\c*sin(\t)});

  % principal axes (hidden inner parts dashed)
  \draw[red,thick,dashed] (0,0,0) -- (\a,0,0);
  \draw[->,red,thick] (\a,0,0) -- (3.6,0,0) node[anchor=west] {$x_1$};
  \draw[green!45!black,thick,dashed] (0,0,0) -- (0,\b,0);
  \draw[->,green!45!black,thick] (0,\b,0) -- (0,2.7,0) node[anchor=south east] {$x_2$};
  \draw[blue!75!black,thick,dashed] (0,0,0) -- (0,0,\c);
  \draw[->,blue!75!black,thick] (0,0,\c) -- (0,0,2.2) node[anchor=south] {$x_3$};
  \fill (0,0,0) circle (1.5pt) node[anchor=north east] {$C$};

  % a mass element and its distances to the x1 and x3 axes
  \coordinate (m) at (1.5,0.2,0.75);
  \draw[orange!85!black,densely dashed] (m) -- (0,0,0.75) node[pos=.5,anchor=north,font=\small] {$r_{\perp3}$};
  \draw[orange!85!black,densely dashed] (m) -- (1.5,0,0) node[pos=.55,anchor=west,font=\small] {$r_{\perp1}$};
  
  \fill[orange!85!black] (m) circle (2pt) node[anchor=south west] {$m_i$};

  \begin{scope}[tdplot_screen_coords]
    \node[anchor=west,align=left] at (3.9,0.2)
      {主轴两两正交，主轴系中\\[2pt]
       $\boldsymbol I=\operatorname{diag}(I_1,I_2,I_3)$\\[4pt]
       $I_k=\sum_i m_i\,r_{\perp k,i}^{2}$\\[4pt]
       {\small 图中椭球沿 $x_1$ 最长、沿 $x_3$ 最短：}\\
       {\small 质量离 $x_1$ 轴最近，$I_1<I_2<I_3$}};
  \end{scope}
\end{tikzpicture}
\end{document}
